% ==============================================================================
% TEMPLATE: BAB KESIMPULAN
% ==============================================================================
% File ini memuat struktur bab Kesimpulan dengan format baku 4-5 poin terurut.
% Setiap butir WAJIB diawali dengan judul poin tebal (\item \textbf{...:}).
% Ganti teks placeholder di bawah dengan kesimpulan substantif pengujian riil.
% ==============================================================================

\chapter{Kesimpulan}
\label{chap:kesimpulan}

Berdasarkan serangkaian kegiatan praktikum mengenai [Topik Praktikum] yang telah dilaksanakan—meliputi implementasi panduan praktikum terbimbing hingga penyelesaian studi kasus mandiri pada tugas praktikum—dapat ditarik beberapa kesimpulan penting sebagai berikut:

\begin{enumerate}
    \item \textbf{Efektivitas Implementasi Metode Utama:} Praktikum ini membuktikan bahwa penerapan metode [Nama Metode / Teknologi] mampu mengolah dan menyelesaikan alur kerja sistem secara terstruktur dan terukur sesuai dengan target capaian modul.
    
    \item \textbf{Pentingnya Tahapan Prapemrosesan dan Konfigurasi Awal:} Kualitas luaran akhir sangat ditentukan oleh ketepatan penyiapan parameter lingkungan dan prapemrosesan data masukan, di mana penyesuaian format yang tepat berhasil meminimalisasi terjadinya galat eksekusi.
    
    \item \textbf{Kinerja dan Keunggulan Komparatif:} Pengujian komparasi luaran menunjukkan bahwa pendekatan [Nama Metode Unggul] memberikan performa yang lebih stabil dan adaptif dibandingkan metode dasar, dibuktikan oleh pencapaian metrik evaluasi yang optimal.
    
    \item \textbf{Resiliensi Sistem dan Mitigasi Kendala:} Hambatan teknis yang ditemukan selama eksperimen (seperti fluktuasi galat atau keterbatasan sumber daya komputasi) berhasil dimitigasi melalui optimasi konfigurasi parameter dan validasi alur bertahap.
    
    \item \textbf{Relevansi Praktis Terhadap Kebutuhan Industri / Riset:} Secara keseluruhan, pemahaman mendalam terhadap ekosistem teknologi ini memberikan fondasi yang kuat bagi mahasiswa dalam merancang dan mengembangkan solusi komputasi yang efisien pada skala permasalahan yang lebih kompleks.
\end{enumerate}
